% nl edition of the Book 2 image credits (translated from image-credits-book2.tex
% by the Book 2 nl agent, 2026-10-25); every credit, name and licence kept exactly.
% Image credits for the photographs and AI illustrations of Book 2 (High School).
% Machine-readable license records live in images/book2/CREDITS.md; keep the
% two files in sync. The Book 2 writer owns this file: add one \item per
% photograph (inside the itemize, by chapter) as photographs land.
\chapter*{Verantwoording van de afbeeldingen}

De tekeningen, kaarten en grafieken in dit boek zijn origineel; kaarten en
grafieken die op gepubliceerde gegevens berusten, noemen hun bron in het
onderschrift of in deze lijst. De foto's worden gebruikt onder hun
respectieve vrije licenties, met dank aan hun makers; de volledige
bronlinks en licentie-URL's van elke foto staan in
\texttt{images/book2/CREDITS.md} in de repository van het boek.

\bigskip
Naast de foto's en de originele tekeningen zijn sommige figuren door AI
gegenereerde illustraties, gemaakt met de beeldgeneratie van OpenAI op
basis van prompts die de redactie van het boek heeft geschreven, en vóór
opname op wetenschappelijke juistheid gecontroleerd. Ze tonen alleen
algemene scènes, nooit een echte plek, een echte persoon of een echt
exemplaar. De volledige prompt van elke gegenereerde afbeelding staat in
\texttt{images/book2/ai/PROMPTS.md} in de repository.

\bigskip
\noindent Foto's en gegevens, per hoofdstuk:
\begin{itemize}\small
  \item Hfst.~1 (de Aarde meten): de Aarde gezien vanuit Apollo~17; foto van NASA, publiek domein (via Wikimedia Commons).
  \item Hfst.~1: een herhalingscirkel, plaat uit J.-B. Delambre, \emph{Base du système métrique décimal}, deel~2 (1807); publiek domein (via Wikimedia Commons, uit BnF Gallica).
  \item Hfst.~1: lengtes van een graad berekend op de ellipsoïde WGS~84 (EPSG-dataset).
  \item Hfst.~2 (een bewoonbare planeet): Venus vanuit Mariner~10 (PIA23791), NASA/JPL-Caltech, publiek domein; Mount Sharp vanaf de rover Curiosity (PIA16769), NASA/JPL-Caltech/MSSS, publiek domein (via Wikimedia Commons).
  \item Hfst.~2: gegevens over de planeten uit de fact sheets van NASA; fasegrenzen van water uit de IAPWS-releases SR1-86 en R14-08; grenzen van de bewoonbare zone uit Kopparapu en anderen (2013); Proxima~b uit het NASA Exoplanet Archive.
  \item Hfst.~3 (het energiebudget van de Aarde): smeltplassen op Arctisch zee-ijs (missie ICESCAPE, 2011); foto van NASA Goddard Space Flight Center, publiek domein (via Wikimedia Commons).
  \item Hfst.~3: waarden van het energiebudget uit het zesde evaluatierapport van het IPCC, werkgroep~I, figuur~7.2; albedo's van de University of Colorado Boulder (MOSAiC) en het National Snow and Ice Data Center; warmtetransport uit Trenberth en Caron (2001); netto primaire productie uit Field en anderen (1998); thermochemische gegevens uit het NIST Chemistry WebBook.
  \item Hfst.~4 (waarom wolken ontstaan): stapelwolken bij mooi weer, foto van Ralph F. Kresge, NOAA; het aambeeld van een cumulonimbus, NOAA National Severe Storms Laboratory; beide publiek domein (via Wikimedia Commons).
  \item Hfst.~4: verzadiging berekend uit de IAPWS-releases; druppelgroottes van het NASA Jet Propulsion Laboratory.
  \item Hfst.~5 (stormen en tropische cyclonen): orkaan Isabel vanuit het Internationaal Ruimtestation, foto van Mike Trenchard, NASA Johnson Space Center; een tornado in Kansas, foto van Sean Waugh, NOAA National Severe Storms Laboratory; bliksem, foto van C.~Clark, NOAA National Severe Storms Laboratory; alle publiek domein (via Wikimedia Commons).
  \item Hfst.~5: banen van cyclonen uit IBTrACS versie~4r01 (Gahtan en anderen, 2024; Knapp en anderen, 2010; NOAA National Centers for Environmental Information); categorieën van orkanen van het NOAA National Hurricane Center; energie van orkanen en voorwaarden voor hun ontstaan van het NOAA Atlantic Oceanographic and Meteorological Laboratory; bliksemfrequentie uit Christian en anderen (2003); gegevens over tornado's van het NOAA Storm Prediction Center.
  \item Hfst.~6 (de gelaagde oceaan): een bemonsteringsrozet, foto van Pablo Clemente-Colon, NOAA Photo Library; een profielboei die vanaf een onderzoeksschip te water wordt gelaten, NOAA Global Ocean Monitoring and Observing Program; chlorofyl in de oceaan 2002--2010, Jesse Allen en Robert Simmon, NASA Earth Observatory; alle publiek domein (via Wikimedia Commons).
  \item Hfst.~6: oceaanprofielen uit de World Ocean Atlas 2023 (Reagan en anderen, 2024; NOAA National Centers for Environmental Information); dichtheid van zeewater uit de internationale toestandsvergelijking (Fofonoff en Millard, 1983); volume van de oceaan uit Charette en Smith (2010); gegevens over de boeien van het Argo Program Office; opwelling van de NOAA National Ocean Service en Kessler (2006).
  \item Hfst.~7 (de bodem, een langzaam gevormde hulpbron): een laterietprofiel op gneis in Kerala, foto van Werner Schellmann, CC~BY-SA~2.5; saproliet in centraal Madagaskar, foto van Diorit, CC~BY-SA~3.0; een bodemprofiel, USDA Natural Resources Conservation Service, publiek domein (alle via Wikimedia Commons).
  \item Hfst.~7: snelheden van bodemvorming en erosie uit Montgomery (2007); mineraalformules uit de Mineralogy Database (webmineral.com); dikte van saproliet uit de USGS Ground Water Atlas of the United States.
  \item Hfst.~8 (stroomgebieden, rivieren en overstromingen): een hydroloog die het debiet meet en een meetstation, U.S. Geological Survey, publiek domein (via Wikimedia Commons).
  \item Hfst.~8: jaarlijkse pieken en dagelijks debiet van de Potomac bij Washington, U.S. Geological Survey National Water Information System, station 01646500; debiet en stroomgebied van de Mississippi uit Kammerer (1990, USGS); debiet van de Amazone uit Callède en anderen (2010); veiligheid bij overstromingswater van de NOAA National Weather Service.
  \item Hfst.~9 (zoetwatervoorraden): het Aralmeer in 2000 en 2018, NASA Earth Observatory (MODIS/Terra); geïrrigeerde velden bij Sublette, Kansas, NASA (ASTER/Terra), publiek domein (via Wikimedia Commons).
  \item Hfst.~9: wereldwijde watervoorraden naar Shiklomanov (1993), U.S. Geological Survey Water Science School; de wereldwijde waterkringloop naar Trenberth en anderen (2007), zoals afgedrukt in Trenberth en Asrar (2014); ouderdommen van het Nubische grondwater uit Sturchio en anderen (2004); veranderingen van de High Plains-aquifer uit McGuire (2017, USGS); de waterbalans van het Aralmeer uit Gaybullaev en anderen (2012) en Micklin (2007).
  \item Hfst.~10 (van continentverschuiving tot oceaanbodemspreiding): het boorschip \emph{JOIDES Resolution}, Ocean Drilling Program, Texas A\&M University, publiek domein (via Wikimedia Commons).
  \item Hfst.~10: de polariteitstijdschaal van Cande en Kent (1995); de boorresultaten van Leg~3 van het Deep Sea Drilling Project (1970); de ouderdom van de oceaanbodem uit Seton en anderen (2020), EarthByte, licentie CC~BY~4.0.
  \item Hfst.~11 (seismische golven onthullen het binnenste): het plaatsen van een tijdelijke seismometer in Puerto Rico; foto van Thomas Pratt, U.S. Geological Survey, publiek domein (via Wikimedia Commons).
  \item Hfst.~11: het Preliminary Reference Earth Model van Dziewonski en Anderson (1981); empirische looptijden van P- en S-golven van de Australian Passive Seismic Server (AuSPASS).
  \item Hfst.~12 (lithosfeer en asthenosfeer): een ijsberg, NOAA, publiek domein (USGS Water Science School).
  \item Hfst.~12: warmtestroom uit Davies en Davies (2010), licentie CC~BY~3.0; warmtestroomprovincies uit Roy, Blackwell en Decker (1972); de mantelsolidus van Hirschmann (2000); dikte en dichtheid van de korst uit Christensen en Mooney (1995); hoogten uit ETOPO1 (NOAA).
  \item Hfst.~13 (platen en hun bewegingen): een satellietplaatsbepalingsstation in de Cascade Range; foto van Dave Sherrod, U.S. Geological Survey, publiek domein (via Wikimedia Commons). Schaduwreliëf van de bodem van de noordelijke Grote Oceaan, National Geophysical Data Center, NOAA, publiek domein (via Wikimedia Commons).
  \item Hfst.~13: rotatiepolen van de platen NNR-MORVEL56 van Argus, Gordon en DeMets (2011); stationsposities van het Nevada Geodetic Laboratory (Blewitt, Hammond en Kreemer, 2018); ouderdommen van vulkanen van Clague en Dalrymple (1987), U.S. Geological Survey; plaatgrenzen van Bird (2003); land van Natural Earth.
  \item Hfst.~14 (oceaanruggen): kussenlava op de Galápagosrift, NOAA Okeanos Explorer Program, publiek domein; een zwarte roker op de Juan de Fucarug, University of Washington en NOAA, publiek domein; gabbro, foto van Darla Sondrol (GeoDIL), CC0; peridotiet in basalt, foto van James St.~John, CC~BY~2.0 (alle via Wikimedia Commons).
  \item Hfst.~14: dieptes van de oceaanbodem uit ETOPO1 (NOAA) en ouderdommen uit Seton en anderen (2020), EarthByte, licentie CC~BY~4.0; dieptewetten van Parsons en Sclater (1977); de mantelsolidus van Hirschmann (2000); rotatiepolen NNR-MORVEL56 van Argus, Gordon en DeMets (2011).
  \item Hfst.~15 (subductiezones): blauwschist en eclogiet, foto's van James St.~John, CC~BY~2.0; tsunamischade in Kesennuma, Japan, 2011, foto van Dylan McCord, U.S. Navy, publiek domein (alle via Wikimedia Commons).
  \item Hfst.~15: aardbevingen uit de aardbevingscatalogus van de USGS (ComCat); plaatgrenzen en convergentiesnelheden van Bird (2003); ouderdommen van de oceaanbodem uit Seton en anderen (2020), EarthByte, licentie CC~BY~4.0; plaatsnelheden NNR-MORVEL56 van Argus, Gordon en DeMets (2011); land van Natural Earth.
  \item Hfst.~16 (mineralen): kwarts, orthoklaas, biotiet en olivijn, foto's van Shannon Heinle, Darla Sondrol, Dexter Perkins en Nessa Eull (Geology Digital Image Library), CC0; basalt in slijpplaatje, foto's van Chmee2, CC~BY-SA~3.0; kwarts in een graniet in slijpplaatje, foto van Michael C.~Rygel, CC~BY-SA~3.0 (alle via Wikimedia Commons).
  \item Hfst.~16: de samenstelling van de continentale korst van Rudnick en Gao (2014); standaardatoommassa's van IUPAC/CIAAW; mineraalformules en dichtheden van webmineral.com; dichtheden van granietachtige gesteenten van de Sierra Nevada van de U.S. Geological Survey (Open-File Report 91-570-A).
  \item Hfst.~17 (stollingsgesteenten): basaltzuilen van de Giant's Causeway, foto van Adam Bishop, CC~BY-SA~3.0; El Capitan, Yosemite, U.S. Geological Survey, CC0; basalt, andesiet, ryoliet, dioriet, graniet, obsidiaan, een porfirisch gesteente en twee slijpplaatjes, foto's van Nessa Eull, Twyla Baker, Joseph H.~Hartman, Mike Davis, Dexter Perkins en Darla Sondrol (Geology Digital Image Library), CC0 (alle via Wikimedia Commons).
  \item Hfst.~17: afkoeling van gangen en sills naar Delaney (1986), U.S. Geological Survey Open-File Report 86-444; samenstellingen van magma's en de kristallisatievolgorde uit de cursusnota's fysische geologie van S.~A.~Nelson en S.~Earle, \emph{Physical Geology} (CC~BY~4.0); de biografie van Bowen uit de geschiedenis van het Geophysical Laboratory van Carnegie Science.
  \item Hfst.~18 (steenkool, olie en gas): fossiele planten in koolstofrijke schalie, bruinkool, bitumineuze steenkool en antraciet, foto's van James St.~John, CC~BY~2.0; een bruinkoollaag in Roemenië, U.S. Geological Survey, publiek domein (alle via Wikimedia Commons).
  \item Hfst.~18: olievenster, steenkoolvorming en vallen uit S.~Earle, \emph{Physical Geology} (CC~BY~4.0); koolstofgehalte van steenkool van de U.S. Energy Information Administration; het petroleumsysteem naar Magoon en Schmoker (U.S. Geological Survey, 2000); de rift- en zoutbekkens van westelijk Centraal-Afrika uit U.S. Geological Survey Fact Sheet 2023-3028.
  \item Hfst.~19 (de continentale korst): de gneis van Acasta, foto van Pedroalexandrade, CC~BY-SA~3.0; zirkoonkorrels van de Jack Hills, beelden van A.~Cavosie en J.~Valley, publiek domein; gneizen en granieten van het Canadees Schild, foto van Chris Light, CC~BY-SA~4.0 (alle via Wikimedia Commons).
  \item Hfst.~19: samenstelling van de korst van Rudnick en Gao (2014); korstdikte van Christensen en Mooney (1995), White, McKenzie en O'Nions (1992) en Karplus en anderen (2019); uitgestrektheid van de continenten van Cogley (1984); ouderdommen van de oceaanbodem uit Seton en anderen (2020), EarthByte, licentie CC~BY~4.0.
  \item Hfst.~20 (gebergtevorming): de Mount Everest vanuit de spaceshuttle Columbia, NASA, publiek domein; geplooide gesteentelagen in de kloof van de Kali Gandaki, foto van Jean-Marie Hullot, CC~BY~2.0; de Glarus-overschuiving, foto van Christian Heine, CC~BY-SA~3.0; kussenlava van de ofioliet van Semail, foto van Lionel Duchoiselle, CC~BY-SA~3.0 (alle via Wikimedia Commons).
  \item Hfst.~20: vroegere posities van India en Lhasa uit de plaatmodellen van Müller en anderen (2019, 2022), berekend met de GPlates Web Service (EarthByte); hooggelegen land uit ETOPO1; aardbevingen uit de catalogus van de USGS; plaatgrenzen van Bird (2003); land van Natural Earth.
  \item Hfst.~21 (metamorfe gesteenten): leisteenklieving onder de microscoop, foto van Geoloshe, CC0; granaat-glimmerschist van Syros, foto van Graeme Churchard, CC~BY~2.0; kyanietschist en kwarts-biotietgneis, foto's van James St.~John, CC~BY~2.0 (alle via Wikimedia Commons).
  \item Hfst.~21: de stabiliteitsvelden van kyaniet, andalusiet en sillimaniet uit Holdaway en Mukhopadhyay (1993).
  \item Hfst.~22 (magmatisme bij subductiezones): de eruptiepluim van Mount St.~Helens, 18~mei 1980, foto van Donald A. Swanson, USGS, publiek domein; granodioriet, andesiet en de Sierra Nevada, foto's van James St.~John, CC~BY~2.0 (alle via Wikimedia Commons).
  \item Hfst.~22: aardbevingen uit de catalogus van de USGS; vulkanen van het Global Volcanism Program van het Smithsonian; plaatgrenzen van Bird (2003); land van Natural Earth; de solidus van droge peridotiet van Hirschmann (2000), met dichtheden uit PREM.
  \item Hfst.~23 (de afbraak van gebergten): de Appalachen vanuit het Internationaal Ruimtestation, astronautenfoto ISS033-E-22378, NASA, publiek domein; kwartskiezelconglomeraat, foto van James St.~John, CC~BY~2.0 (beide via Wikimedia Commons).
  \item Hfst.~23: denudatiesnelheden van Portenga en Bierman (2011); de sedimentvrachten van de Mississippi van Meade en Moody (2010); de vervalconstante van Ahnert (1970); de Varden Hills, de verzonnen streek van de reeks.
  \item Hfst.~24 (de inwendige warmte van de Aarde): de geothermische centrale van Nesjavellir, IJsland, foto van Gretar Ívarsson, bewerkt door Fir0002, publiek domein (via Wikimedia Commons).
  \item Hfst.~24: warmtestroom van Davies en Davies (2010); radiogene warmte uit de geoneutrinometing van KamLAND (2011) en de huidige warmteafgifte van de vier isotopen; warmteproductie van gesteenten uit Vilà en Fernández (2010).
  \item Hfst.~25 (radiometrische datering en de ouderdom van de Aarde): William Thomson, Lord Kelvin, 1906, foto van Russell and Sons, Londen, Smithsonian Institution Libraries, geen bekende auteursrechtelijke beperkingen; de ijzermeteoriet van Canyon Diablo, foto van Geoffrey Notkin, CC~BY-SA~2.5 (beide via Wikimedia Commons).
  \item Hfst.~25: loodisotopen van meteorieten van Patterson (1956); de door IUGS en IUPAC (2022) aanbevolen vervalconstante van rubidium-87; de geschiedenis van de schattingen naar Holmes (1913, 1927), Patterson (1956) en de U.S. Geological Survey; de Varden Hills, de verzonnen streek van de reeks.
  \item Hfst.~26 (de jonge Aarde en haar atmosfeer): gebande ijzerformatie, Dales Gorge, foto van Graeme Churchard, CC~BY~2.0; stromatolieten in Shark Bay, foto van Paul Harrison, CC~BY-SA~3.0; rode lagen bij Sedona, foto van Wendy from Pennsylvania, CC~BY-SA~2.0 (alle via Wikimedia Commons).
  \item Hfst.~26: de lichtkracht van de jonge Zon van Gough (1981); grenzen en kleuren van de eonen uit de International Chronostratigraphic Chart.
  \item Hfst.~27 (klimaten van het verleden lezen): luchtbelletjes in Antarctisch ijs, foto CSIRO Atmospheric Research, CC~BY~3.0; planktonische foraminiferen, foto van Marci Robinson, USGS, publiek domein; pollenkorrels, Dartmouth Electron Microscope Facility, publiek domein (alle via Wikimedia Commons).
  \item Hfst.~27: de temperatuur van EPICA Dome~C (Jouzel en anderen, 2007) en het koolstofdioxide van de Antarctische ijskernen (Bereiter en anderen, 2015) via NOAA Paleoclimatology; de LR04-stapel (Lisiecki en Raymo, 2005); het reliëf van ETOPO1 (NOAA); NASA GISS (GISTEMP v4); referentieperioden uit het zesde evaluatierapport van het IPCC, werkgroep~I, hoofdstuk~2.
  \item Hfst.~28 (klimaatverandering): het Mauna Loa Observatory; foto van Brian Vasel, NOAA, publiek domein (via Wikimedia Commons).
  \item Hfst.~28: Charles David Keeling en John Chin op Mauna Loa, 1997; foto van John Miller, NOAA, publiek domein (via Wikimedia Commons).
  \item Hfst.~28: gegevens van het NOAA Global Monitoring Laboratory en de Scripps Institution of Oceanography (koolstofdioxide), NASA GISS (GISTEMP v4), het NOAA Laboratory for Satellite Altimetry (zeespiegel), Bereiter en anderen (2015) via NOAA Paleoclimatology (ijskernen), en het zesde evaluatierapport van het IPCC, werkgroep~I (forceringen, toeschrijving, projecties).
\end{itemize}
